% Generated draft: requires mathematical and source review.
\documentclass[11pt]{article}
\usepackage[margin=25mm]{geometry}
\usepackage{fontspec}
\setmainfont{Latin Modern Roman}
\usepackage{amsmath,amssymb,mathtools,unicode-math}
\setmathfont{Latin Modern Math}
\usepackage{enumitem,needspace,xurl,hyperref}
\setlength{\parindent}{0pt}
\setlength{\parskip}{5pt}
\setlength{\emergencystretch}{3em}
\newcommand{\N}{\mathbb{N}}
\newcommand{\Z}{\mathbb{Z}}
\newcommand{\Q}{\mathbb{Q}}
\newcommand{\R}{\mathbb{R}}
\newcommand{\C}{\mathbb{C}}
\newcommand{\F}{\mathbb{F}}
\newcommand{\A}{\mathbb{A}}
\newcommand{\PP}{\mathbb{P}}
\newcommand{\E}{\mathbb{E}}
\newcommand{\eps}{\varepsilon}
\newcommand{\dd}{\,\mathrm{d}}
\begin{document}
\section*{UnsolvedMath ID: 1100309}
\textbf{Source problem code:} catalogue entry
\subsection{Definitions and mathematical statement}
For an integer $n\geq2$, let $F_n=\langle a_1,\ldots,a_n\rangle$ be the free group of rank $n$. An ordered tuple $(x_1,\ldots,x_n)$ is a basis when the homomorphism $F_n\to F_n$ sending $a_i$ to $x_i$ is an automorphism. Put
\[
B_n=\{(x_1,\ldots,x_n)\in F_n^n:(x_1,\ldots,x_n)\text{ is a basis of }F_n\}.
\]
Use the first-order language of groups $(\cdot,{}^{-1},1)$. The set $B_n$ is definable without parameters if there is a formula $\varphi(x_1,\ldots,x_n)$ in this language whose realizations in $F_n$ are exactly $B_n$. It is definable with parameters if there are a finite tuple $c\in F_n^m$ and a formula $\varphi(x_1,\ldots,x_n;y_1,\ldots,y_m)$ with
\[
B_n=\{x\in F_n^n:F_n\models\varphi(x;c)\}.
\]
\textbf{Source conjecture.} For each integer $n>1$, $B_n$ is not definable. The displayed source sentence does not say whether parameters are permitted, so its two interpretations must be distinguished.

The supplied research summary reports: without parameters, $B_n$ is not definable for every $n\geq2$; with parameters, $B_2$ is definable whereas $B_n$ is not definable for every $n\geq3$. For rank two it attributes parameter-definability to Nielsen's criterion that $[g,h]$ be conjugate to $[a_1,a_2]$ or its inverse, where $[u,v]=u^{-1}v^{-1}uv$. Thus the two interpretations have different answers at rank two. These are reported literature conclusions, not independently verified here.
\subsection{Short English statement}
Can first-order group formulas describe precisely the ordered free bases of $F_n$? The supplied resolution distinguishes formulas with parameters from formulas without them, with a rank-two exception only in the former case.
\subsection{Sources}
\begin{itemize}
\item M. Bestvina, \emph{Questions in Geometric Group Theory} (2004), Question 3.9, PDF page 11: \url{https://www.math.utah.edu/~bestvina/eprints/questions-updated.pdf}; manual transcription, accessed 2026-07-29.
\item Supplied research-summary attributions: Nielsen (1917), rank-two commutator criterion; Pillay, characterization of the generic type; Sklinos, bigness of the non-primitives; Bestvina--Feighn, negligibility or co-negligibility of definable sets and the higher-rank primitive-set obstruction. Apart from Nielsen's year, no publication identifiers are supplied.
\item Local registry \texttt{data/unsolvedmath/problems.json}, permanent ID $1100309$, source code catalogue entry, \emph{Questions in Geometric Group Theory --- Q 3.9}. The record labels the question solved and SOLVED-IN-LITERATURE; the source-list background separately retains NEEDS\_REVIEW status and rights qualifications.
\end{itemize}
\end{document}
